\begin{problem}[第37届全国中学生物理竞赛决赛][导体球间电场增强现象]{126}
    当电场中两个导体球靠近时，导体球之间的电场将明显增强。本题试探讨这一现象。已知真空介电常量为$\varepsilon _ { 0 }$。
    \begin{subproblem}
        设一个半径为$R _ { 0 }$的孤立导体球的球心与一个静止点电荷$Q$相距为$a \ (a > R _ { 0 } )$，求镜像电荷的电量及其位置。
    \end{subproblem}
    \begin{subproblem}
        设导体球置于大小为$E _ { 0 }$的匀强外电场中，该外加电场可看作是由两个相距很远的等量异号点电荷$\pm Q$在其连线中点处产生的。试证导体球的感应电荷的作用等效于两个镜像电荷形成的电偶极子，并求该电偶极子的偶极矩与外场$E _ { 0 }$之间的关系。
    \end{subproblem}
    \begin{subproblem}
        设导体球外两等量异号点电荷$\pm q$的间距为$\Delta l$，它们形成的电偶极子$q \Delta l$沿径向放置在半径为$R _ { 0 }$的导体球附近，偶极子中心与导体球中心相距$r \ ( r >> \Delta l )$，求该偶极子在导体球中镜像电偶极矩的大小。
    \end{subproblem}
    \begin{subproblem}
        在大小为$E _ { 0 }$的匀强外电场中，沿外电场方向放置两个半径为$R _ { 0 }$的导体球，两球心相距$r \ ( r > 2 R _ { 0 } )$。求导体球外过两球心的平面内任一点$P'(x, y)$处的电势分布和在两球心连线方向的场强分布（可用递推式表示），取两球心连线中点为坐标原点，连线方向为$X$轴。
    \end{subproblem}
    \begin{subproblem}
        试证：在不考虑击穿放电的情形下，当上述两导体缓慢无限靠近时，两球连心线中点的电场会因静电感应而趋于无限大。
    \end{subproblem}
\end{problem}









